% "Geometric Dark Energy: Phenomenological Viability,
%  Systematic Closures, and Requirements for Completion"
%
% Houston Golden — 2026

\documentclass[aps,prd,twocolumn,superscriptaddress,showpacs,preprintnumbers,
  nofootinbib,longbibliography,floatfix]{revtex4-2}

\usepackage{amsmath,amssymb,amsfonts}
\usepackage{graphicx}
\usepackage{bm}
\usepackage{hyperref}
\usepackage{xcolor}
\usepackage{booktabs}
\usepackage{multirow}
\usepackage{dcolumn}
\usepackage{enumitem}
\usepackage{mathtools}
\usepackage[utf8]{inputenc}
\usepackage{float}
\usepackage{slashed}

\hypersetup{
  colorlinks=true,
  linkcolor=blue!60!black,
  citecolor=blue!60!black,
  urlcolor=blue!60!black
}

% Custom commands
\newcommand{\Leff}{\Lambda_{\rm eff}}
\newcommand{\Lobs}{\Lambda_{\rm obs}}
\newcommand{\Dinf}{\mathcal{D}_{\rm inf}}
\newcommand{\rhocrit}{\rho_{\rm crit}}
\newcommand{\rhoPl}{\rho_{\rm Pl}}
\newcommand{\MPl}{M_{\rm Pl}}
\newcommand{\Neff}{N_{\rm eff}}
\newcommand{\BI}{\gamma}
\newcommand{\Geff}{G_{\rm eff}}
\newcommand{\Gsp}{G_{\rm SP}}
\newcommand{\Veff}{V_{\rm eff}}
\newcommand{\ie}{i.e.}
\newcommand{\eg}{e.g.}
\newcommand{\cf}{cf.}
\newcommand{\etal}{\textit{et~al.}}
\begin{document}

\preprint{HUBIFY-2026-002}

\title{Geometric Dark Energy: Phenomenological Viability,\\
Systematic Closures, and Requirements for Completion}

\author{Houston Golden}
\email{houston@hubify.com}
\affiliation{Independent Researcher, Los Angeles, California, USA}

\date{\today}

\begin{abstract}
We present a phenomenological dark-energy framework rooted in
Einstein-Cartan-Holst (ECH) gravity, report the systematic closure of
four minimal and one near-minimal first-principles route to deriving the
late-time vacuum term, and extract the structural requirements that any
viable geometric dark-energy completion must satisfy.

The minimal ECH framework---with algebraic torsion, fixed Barbero-Immirzi
parameter $\BI = 0.274$, and minimal Dirac coupling---provides a
consistent phenomenological ansatz that reduces fine-tuning from
$10^{120}$ to $\sim\!10^5$ and yields MCMC fits partially reducing the
$H_0$ and $\sigma_8/S_8$ tensions. However, four systematic attempts to
derive the late-time $w = -1$ behavior within the minimal model have all
failed, sharing a common structural origin: algebraic torsion washes out
after elimination, leaving no infrared fingerprint. Additionally, the
most natural next-generation extension---propagating torsion in
Poincar\'e gauge theory (PGT)---reveals a new obstruction: in the
ghost-free $0^-$ axial mode, the mass and matter coupling are locked
($m \propto g_{\rm eff} \propto 1/\sqrt{|t_3|}$), so making the mode
cosmologically light simultaneously drives its coupling $\sim\!10^{29}$
times below gravitational strength, rendering it observationally inert.
No symmetry within PGT protects this mass.

From these closures we extract six structural lessons and formulate
requirements that constrain the space of viable geometric dark-energy
models: the theory must survive reduction, address scale naturalness,
offer a distinctive observable, and---crucially---avoid the
mass-coupling lock in which a light-mass limit automatically decouples
the field from matter.

The framework's phenomenological viability (cosmological fits, scaling
argument, fine-tuning reduction) survives all closures intact. Its
scientific contribution is the combination of a well-tested
phenomenological model with a systematic map of the theory landscape
showing what does not work and why, narrowing the path for any future
geometric completion.
\end{abstract}

\pacs{98.80.-k, 04.50.Kd, 04.60.Pp, 95.36.+x}
\keywords{dark energy, Einstein-Cartan theory, spin-torsion coupling,
cosmological constant problem, Barbero-Immirzi parameter,
negative results, geometric dark energy}

\maketitle
\tableofcontents


%% ============================================================
%%  PART I: FRAMEWORK AND PHENOMENOLOGY
%% ============================================================

%% ============================================================
%%  1. INTRODUCTION
%% ============================================================
\section{Introduction}\label{sec:intro}

The cosmological constant problem---why the observed vacuum energy density
$\rho_\Lambda \sim (2.3\;\text{meV})^4$ is $\sim 10^{120}$ times smaller
than the naive quantum field theory estimate $\sim \MPl^4$~\cite{Weinberg1989}%
---remains one of the deepest unsolved problems in fundamental physics.
Growing tensions within $\Lambda$CDM, including the ${\sim}4.9\sigma$
Hubble tension~\cite{Riess2022,Planck2018params} and $2$--$3\sigma$
$\sigma_8/S_8$ tension~\cite{KiDS2021,DES2024}, together with hints of
dynamical dark energy from DESI~DR2 BAO
measurements~\cite{DESI2024,DESI2025DR2}, motivate exploration of
geometric approaches to dark energy.

Einstein-Cartan-Holst (ECH) gravity---Einstein-Cartan theory supplemented
by the Holst term with Barbero-Immirzi parameter $\BI$---is a natural
starting point for geometric dark energy. The theory is well-motivated by
loop quantum gravity~\cite{Ashtekar2011}, introduces spin-torsion coupling
with parity-violating structure~\cite{Freidel2005,Mercuri2006}, and
provides a non-singular quantum bounce~\cite{Poplawski2012}. The question
is whether this theoretical structure can generate a natural origin for
late-time dark energy.

This paper reports three results. First, we present a phenomenological
dark-energy framework built on ECH gravity that is observationally
consistent and reduces fine-tuning significantly, with MCMC fits partially
reducing the $H_0$ and $\sigma_8$ tensions. Second, we report the
systematic failure of four minimal first-principles routes to derive
the late-time $w = -1$ behavior. Third, we report an assessment
of the most natural next-generation extension---propagating torsion in
Poincar\'e gauge theory (PGT)---finding that it is structurally viable but
phenomenologically empty for dark energy due to a \textit{mass-coupling
lock}: the same parameter that makes the torsion mode cosmologically
light also suppresses its matter coupling to $\sim\!10^{29}$ times below
gravitational strength, rendering it observationally inert. From these
results we extract six structural lessons and formulate concrete
requirements for any viable geometric dark-energy completion.

The combination establishes the current scientific status of the program:
\begin{itemize}[leftmargin=*]
  \item The minimal ECH framework is a \textit{motivated phenomenological
    model}, not a first-principles derivation of dark energy.
  \item The minimal model and its most natural extension have both been
    tested and found insufficient, for distinct structural reasons.
  \item The closures map the theory landscape and identify nontrivial
    requirements---including avoidance of the mass-coupling lock---that
    constrain future geometric dark-energy proposals.
\end{itemize}

\subsection{Relationship to prior work}

This paper supersedes an earlier version~\cite{Golden2026v1} that
presented the same phenomenological framework but was organized around
the minimal ECH model as if it were a sufficient foundation. The present
version incorporates the results of a structured derivation program
(documented in a companion technical note~\cite{Golden2026supplement})
that tested four minimal routes and found all to be negative. The
phenomenological results are unchanged; the framing, interpretation, and
forward-looking directions are substantially revised.


%% ============================================================
%%  2. EINSTEIN-CARTAN-HOLST FRAMEWORK
%% ============================================================
\section{Einstein-Cartan-Holst Framework}\label{sec:ech}

\subsection{Action and field content}\label{sec:action}

The starting point is the Einstein-Cartan-Holst action minimally coupled
to $N_f$ Dirac fermions:
\begin{align}\label{eq:ECH}
  S_{\rm ECH} &= \frac{\MPl^2}{2}\int
    \bigl[\varepsilon_{IJKL} + \tfrac{2}{\BI}\,\eta_{I[K}\eta_{L]J}\bigr]\,
    e^I\wedge e^J\wedge F^{KL}[\omega] \nonumber\\
  &\quad + \int d^4x\,|e|\,\bigl[\bar\psi\,i\gamma^\mu D_\mu\psi
    - m\bar\psi\psi\bigr],
\end{align}
with $D_\mu = \partial_\mu + \tfrac{1}{4}\omega_\mu^{IJ}\gamma_I\gamma_J$
the Lorentz-covariant derivative and $\BI = 0.274 \pm 0.020$ the
Barbero-Immirzi parameter fixed by loop quantum gravity~\cite{Ashtekar2011}.

The first term in Eq.~\eqref{eq:ECH} is the Holst action: the Palatini
(first-order) formulation of general relativity supplemented by the Holst
term $(1/\BI)\varepsilon_{IJKL}e^I e^J F^{KL}$, which is parity-odd and
classically equivalent to a topological term (the Nieh-Yan density) on
shell. The second term is the standard minimal Dirac coupling.

\subsection{Torsion equation of motion}\label{sec:torsion}

In ECH gravity, the spin connection $\omega^{IJ}_\mu$ is an independent
dynamical variable, and its equation of motion is algebraic---torsion is
\textit{not propagating}. Solving the torsion equation exactly
decomposes the connection as $\omega = \mathring\omega + K$, where
$\mathring\omega$ is the torsion-free Levi-Civita connection and $K$ is
the contorsion tensor determined algebraically by the fermion spin
density~\cite{Freidel2005,Shapiro2002}.

\subsection{Reduced action after torsion elimination}\label{sec:reduced}

Substituting the algebraic torsion solution back into the action yields
the reduced action~\cite{Freidel2005,ShapiroTeixeira2014}:
\begin{equation}\label{eq:Sreduced}
  S_{\rm red} = S_{\rm EH}[g] + S_{\rm Dirac}[g,\psi]
  + S_{\rm 4f}[\psi;\BI],
\end{equation}
where $S_{\rm EH}$ is the standard Einstein-Hilbert action,
$S_{\rm Dirac}$ is the standard minimally-coupled Dirac action with
the Levi-Civita connection, and $S_{\rm 4f}$ is the torsion-induced
four-fermion contact interaction:
\begin{equation}\label{eq:L4f}
  \mathcal{L}_{\rm 4f}
  = -\Geff\Bigl(\bar\psi\gamma^\mu\gamma^5\psi
    + \tfrac{1}{\BI}\,\bar\psi\gamma^\mu\psi\Bigr)^{\!2},
  \quad
  \Geff = \frac{3\kappa^2}{16}\,\frac{\BI^2}{\BI^2+1}.
\end{equation}
This perfect-square structure~\cite{Golden2026supplement} constrains the
Fierz channel decomposition. At $\BI = 0.274$, the vector admixture
$1/\BI \approx 3.65$ dominates; the scalar/pseudoscalar channel coupling
$\Gsp \propto (\BI^2 - 1)/(\BI^2 + 1)$ is repulsive (negative).

\paragraph{Critical structural observation.}
After torsion elimination, \textit{all} $\BI$-dependence resides in the
four-fermion term $S_{\rm 4f}$. The Dirac kinetic operator
$i\gamma^\mu\mathring\nabla_\mu$ is the standard Levi-Civita operator,
independent of $\BI$. This structural fact is the root cause of the
failures reported in Sec.~\ref{sec:closures}: the Holst sector's content
does not survive torsion elimination in any form that can generate
late-time vacuum physics at one loop or below.


%% ============================================================
%%  3. PHENOMENOLOGICAL DARK-ENERGY ANSATZ
%% ============================================================
\section{Phenomenological Dark-Energy Ansatz}\label{sec:ansatz}

\subsection{The scaling argument}\label{sec:scaling}

The ECH framework contains a parity-odd vacuum operator
$\varepsilon^{abcd}K_{ab}R_{cd}$ sourced by the fermion spin density via
torsion. At early times (Planck-density bounce), this operator has a
natural magnitude set by $\alpha/M \sim \mathcal{O}(\kappa^2)$, where
$\alpha/M$ is the parity-odd coefficient in the effective Lagrangian.

Inflationary expansion after the bounce dilutes this operator by a factor
$\Dinf \sim e^{-3N_{\rm tot}}$, where $N_{\rm tot} \sim 92$ $e$-folds.
The resulting late-time dark-energy density is parametrically
\begin{equation}\label{eq:rhoL}
  \rho_\Lambda \approx \Xi\,\MPl^4, \qquad
  \Xi \equiv \bigl[(\alpha/M)\,\MPl\bigr]\,\Dinf,
\end{equation}
giving $\Xi \sim 10^{-122}$ for appropriate values of $\alpha/M$ and
$N_{\rm tot}$.

\paragraph{What this is.}
Equation~\eqref{eq:rhoL} is a \textit{phenomenological scaling ansatz}.
It traces a dimensional chain from the Planck scale through inflationary
dilution to the observed vacuum energy. It reduces fine-tuning from
$10^{120}$ (the cosmological constant problem in $\Lambda$CDM) to
$\sim\!10^5$ (sensitivity to $N_{\rm tot}$ within $\Delta N \approx 4$
$e$-folds).

\paragraph{What this is not.}
Equation~\eqref{eq:rhoL} is \textit{not} a derivation of $w = -1$.
The parity-odd operator is sourced by the spin density, which vanishes
at late times ($K_{ab} \to 0$). For the residual to persist as a true
vacuum energy, one must show that integrating out the spin degrees of
freedom generates an IR-constant term in the effective action. Four
attempts to establish this have failed (Sec.~\ref{sec:closures}). We
therefore treat $w = -1$ as a phenomenological assumption, consistent
with observations ($|w+1| \lesssim 10^{-2}$ from DESI+Planck).

\subsection{Cosmological parameters}\label{sec:cosmo}

The phenomenological framework motivates a $\Lambda$CDM$+\Delta\Neff$
extension: the bounce scenario qualitatively predicts additional
relativistic species from particle production near the Planck density,
but quantitative estimates are not yet available. The extension adds
one effective parameter ($\Delta\Neff$) to $\Lambda$CDM; a second
parameter ($(\omega/H)_0$) is bounded by CMB isotropy to
$< 5\times 10^{-11}$~\cite{Saadeh2016} and contributes negligibly.

\subsubsection{Datasets}

We analyze four dataset combinations of increasing constraining power:
\begin{enumerate}[leftmargin=*]
\item \textit{Planck-only}: Planck~2018 NPIPE, specifically
  \texttt{plikHM\_TTTEEE} + \texttt{lowl} + \texttt{lowE} +
  \texttt{lensing}~\cite{Planck2018params}.
\item \textit{Planck + BAO}: Adding DESI~2024 DR1 (seven redshift bins,
  $z_{\rm eff} = 0.30$--$2.33$)~\cite{DESI2024}, supplemented by 6dFGS
  ($z_{\rm eff}=0.106$) and SDSS MGS ($z_{\rm eff}=0.15$).
\item \textit{Planck + BAO + SN}: Adding Pantheon+ (1701 light curves,
  1550 unique SNe, $0.001 < z < 2.26$).
\item \textit{Tension dataset}: Adding SH0ES $H_0$
  prior~\cite{Riess2022} + DES~Y3 $S_8$~\cite{DES2024}.
\end{enumerate}

\subsubsection{MCMC configuration}

Parameter estimation uses Cobaya~\cite{Cobaya2021} (v3.5 original,
v3.6.1 verification) with CAMB, convergence criterion $\hat{R}-1<0.01$,
and $>1000$ effective samples per parameter. The theory engine is
unmodified stock CAMB; curvature is fixed to $\Omega_k = 0$.
The MCMC analysis is phenomenologically equivalent to any
$\Lambda$CDM$+\Delta\Neff$ extension---the tension reduction is a
property of the parameter space, not uniquely of spin-torsion cosmology.

\subsubsection{Results}

MCMC fits to the full tension dataset yield:
\begin{align}
  H_0 &= 69.2 \pm 0.8 \;\text{km}\;\text{s}^{-1}\;\text{Mpc}^{-1},
    \label{eq:H0}\\
  \sigma_8 &= 0.785 \pm 0.016, \label{eq:s8}\\
  S_8 &= 0.80 \pm 0.02. \label{eq:S8}
\end{align}
These partially reduce the $H_0$ tension (from $4.9\sigma$ to
$2.9\sigma$ vs.\ SH0ES) and the $S_8$ tension (from $\sim\!3\sigma$ to
$\sim\!1\sigma$ vs.\ weak lensing).

\subsubsection{Independent verification}\label{sec:verification}

Independent verification using Cobaya~v3.6.1 with Planck NPIPE CamSpec
TTTEEE + lowl TT/EE + lensing has produced two frozen dataset
combinations with publication-quality convergence:

\begin{table}[h]
\caption{\label{tab:verification} Independent verification results
(Cobaya~v3.6.1, CAMB~v1.6.5). Posterior means $\pm\,1\sigma$.}
\begin{ruledtabular}
\begin{tabular}{lcc}
\textbf{Parameter} & \textbf{Full-tension} & \textbf{Planck+BAO+SN} \\
\hline
$H_0$ [km/s/Mpc] & $67.68 \pm 1.06$ & $67.79 \pm 1.09$ \\
$\Delta\Neff$ & $-0.020 \pm 0.169$ & $+0.065 \pm 0.17$ \\
$\sigma_8$ & $0.803 \pm 0.008$ & $0.812 \pm 0.009$ \\
$S_8$ & $0.814 \pm 0.008$ & $0.831 \pm 0.018$ \\
$\Omega_m$ & $0.308 \pm 0.005$ & $0.312 \pm 0.006$ \\
\hline
Chains & 6 & 6 \\
Total samples & 176{,}840 & 132{,}949 \\
Worst $\hat{R}-1$ & 0.001 & 0.003 \\
Min ESS & 4{,}744 & 4{,}692 \\
\end{tabular}
\end{ruledtabular}
\end{table}

Both datasets find $\Delta\Neff$ consistent with zero. The $H_0$ values
($67.68$ and $67.79$) are consistent with Planck $\Lambda$CDM
($67.36 \pm 0.54$) at $0.3\sigma$, confirming that the $\Delta\Neff$
extension alone does not resolve the Hubble tension---the tension
reduction in the original analysis (Eqs.~\ref{eq:H0}--\ref{eq:S8}) is
driven by the SH0ES prior.
The difference between the original central value ($H_0 = 69.2$,
Eqs.~\ref{eq:H0}--\ref{eq:S8}) and the verification
($H_0 = 67.68$, Table~\ref{tab:verification}) reflects different
Planck likelihood implementations (\texttt{plikHM} vs.\ \texttt{CamSpec})
and dataset vintages; both values lie within their respective
$1\sigma$ uncertainties and agree at $1.1\sigma$.

\textit{Comparison with independent EC torsion analyses.}---Liu
\etal~\cite{ECTorsionDESI2025} constrained an EC torsion model using
DESI~DR2 + PantheonPlus + DES~Y5 + Planck, finding $H_0 = 68.41 \pm
0.32$ with torsion preferred over $\Lambda$CDM by AIC
($\Delta\text{AIC} = -5.7$ to $-6.6$). Our verification agrees at
$0.5\sigma$ in $H_0$ and $0.4\sigma$ in $\sigma_8$, providing
independent cross-validation of the EC cosmological framework.

\subsubsection{Bayesian model comparison}

\begin{table}[h]
\caption{\label{tab:modelcomp} Bayesian model comparison (effective
$\chi^2_{\rm eff} \equiv -2\ln\mathcal{L}_{\rm max}$ for the full
tension dataset).}
\begin{ruledtabular}
\begin{tabular}{lccc}
Model & $k$ & $\chi^2_{\rm eff}$ & $\ln B$ \\
\hline
$\Lambda$CDM & 6 & 1156.2 & 0.0 \\
$w$CDM & 7 & 1154.8 & $-0.5$ \\
\textbf{ECH phenom.} & \textbf{7} & \textbf{1148.3} & $\mathbf{+4.8}$ \\
\end{tabular}
\end{ruledtabular}
\end{table}

The Bayes factor is dataset-dependent: $\ln B = -1.2$ (Planck+BAO), $+2.1$
(Planck+BAO+SH0ES), $+4.8$ (full tension). The preference emerges
specifically when tension data are included, because the extra parameter
$\Delta\Neff$ partially reduces these tensions. With Planck+BAO alone,
the Occam penalty slightly disfavors the extension. This dataset
dependence is expected and not a weakness: it reflects the model's
statistical flexibility rather than unique predictive power.

\subsection{Observational consistency checks}\label{sec:obs}

The framework accommodates several observational features as consistency
checks---not predictions.

\subsubsection{Cosmic birefringence}

The framework's parity-odd structure is qualitatively consistent with
the $2.4$--$2.9\sigma$ cosmic birefringence detections from
Planck~\cite{Eskilt2022} and ACT~DR6~\cite{DiegoPalazuelos2025}
(combined: $\beta = 0.242^\circ \pm 0.061^\circ$, $3.9\sigma$).
However, the minimal ECH framework has \textit{no derived photon-torsion
coupling}: matching the observed $\beta$ requires an assumed coupling
with $\mathcal{O}(1)$ coefficient, generic to any ALP model (see
Route~S1 closure, Sec.~\ref{sec:S1}). This is a consistency check,
not a prediction.

\subsubsection{Galaxy spin asymmetry}

A galaxy spin dipole asymmetry ($A_0 \sim 0.003$) is observed
empirically, but the CMB isotropy bound constrains the framework's
rotation-driven amplitude to $A_0 \lesssim 10^{-12}$---an
order-of-magnitude gap that prevents any first-principles connection.
The signal is empirical; its relation to the framework remains
qualitative at best.

\subsubsection{Fine-tuning reduction}

A Monte Carlo sensitivity scan ($10^5$ samples) confirms that the
scaling ansatz Eq.~\eqref{eq:rhoL} reduces fine-tuning from $10^{120}$
to $\sim\!10^5$: the viable range $N_{\rm tot} \in [79,\,95]$ produces
$\rho_\Lambda$ within two orders of magnitude of the observed value,
with $N_{\rm tot}$ as the single controlling parameter ($|\rho_s| =
0.996$). The residual tuning is concentrated in a single initial
condition with potential dynamical answers through the
bounce-to-inflation transition.

\begin{table}[h]
\caption{\label{tab:finetuning} Fine-tuning comparison.}
\begin{ruledtabular}
\begin{tabular}{lcc}
Model & Natural Scale & FT Score \\
\hline
$\Lambda$CDM & $\MPl^4$ & $\sim 10^{120}$ \\
Quintessence & $M_{\rm EW}^4$ & $\sim 10^{60}$ \\
$f(R)$ Gravity & $H_0^4$ & $\sim 10^{40}$ \\
\textbf{ECH phenom.} & $(\alpha/M)\Dinf\,\MPl^4$\,\textsuperscript{$\dagger$} & $\sim\mathbf{10^{5}}$ \\
\end{tabular}
\end{ruledtabular}
\textsuperscript{$\dagger$}The ECH natural scale is a parametric
estimate incorporating the scaling ansatz~\eqref{eq:rhoL}, not an
independently derived scale; the fine-tuning score reflects sensitivity
to $N_{\rm tot}$ within ${\sim}4$ $e$-folds.
\end{table}

\subsubsection{Limit behavior}

Internal consistency is verified: $T^{abc}\to 0$ recovers GR,
$\Dinf\to 1$ reproduces the standard CC problem, and $\BI\to\infty$
recovers ECSK theory.


%% ============================================================
%%  PART II: SYSTEMATIC MODEL TESTING
%% ============================================================

%% ============================================================
%%  4. THE DERIVATION PROGRAM
%% ============================================================
\section{The Derivation Program}\label{sec:program}

The central question is whether late-time $w = -1$ can be derived from
first principles within the minimal ECH$+$Dirac model. We tested four
routes, each with predefined gates and kill criteria frozen before
computation~\cite{Golden2026supplement}.

\subsection{Four routes tested}\label{sec:routes}

\begin{description}[style=nextline,leftmargin=1em]
  \item[Track B (NJL condensate)]
    Test whether a pseudoscalar fermion condensate
    $\langle\bar\psi i\gamma^5\psi\rangle$ forms via the Nambu-Jona-Lasinio
    mechanism from the torsion-induced four-fermion interaction and persists
    as late-time vacuum energy.

  \item[Branch G v1 (strict one-loop effective action)]
    Compute the one-loop gravitational effective action $\Gamma[e]$ after
    torsion elimination and fermion integration, and test for
    $\BI$-dependent vacuum content beyond standard cosmological-constant
    renormalization.

  \item[Route T1 (dynamical Immirzi field)]
    Promote $\BI \to \theta(x)$ to a dynamical pseudoscalar and assess
    whether the reduced action contains novel gravitational content not
    reproducible by adding a standard axion-like particle (ALP) to GR.

  \item[Route S1 (parity-sensitive CMB phenomenology)]
    Assess whether the framework's parity-odd operator maps to distinctive,
    testable birefringence observables distinguishable from generic ALP
    predictions.
\end{description}

All four routes were closed with clean negative results. The detailed
closure documentation is in the companion technical
note~\cite{Golden2026supplement}. Section~\ref{sec:closures} summarizes
the results and Sec.~\ref{sec:lessons} extracts the structural lessons.


%% ============================================================
%%  5. CLOSURE RESULTS
%% ============================================================
\section{Closure Results}\label{sec:closures}

\subsection{Track B: condensate route}\label{sec:trackB}

Fierz rearrangement of the four-fermion
interaction~\eqref{eq:L4f} decomposes it into standard NJL channels.
The scalar/pseudoscalar coupling is
\begin{equation}\label{eq:GSP}
  \Gsp = \frac{3\kappa^2}{16}\,\frac{\BI^2 - 1}{\BI^2 + 1},
\end{equation}
which is \textit{repulsive} (negative) for $\BI < 1$. At $\BI = 0.274$,
the NJL gap equation has no nontrivial solution. Even for $\BI > 1$
(attractive sign), the coupling is $\sim\!175\times$ below the NJL
critical coupling, and gravitational catalysis is exponentially
suppressed ($M_* \sim e^{-2100}\,\MPl$). Track~B closes at Gate~1:
no condensate forms.

\subsection{Branch G v1: strict one-loop}\label{sec:branchG}

After torsion elimination, the Dirac operator is
$\slashed{D} = i\gamma^\mu\mathring\nabla_\mu$---the standard
$\BI$-independent Levi-Civita operator. The four-fermion term is quartic
in $\psi$ and does not contribute at one loop. The Seeley-DeWitt
coefficients $a_0, a_1, a_2$ are the standard textbook values with no
$\BI$-dependence~\cite{Vassilevich2003}. Branch~G v1 closes: no novel
Holst content at strict one-loop.

\subsection{Route T1: dynamical Immirzi field}\label{sec:T1}

Promoting $\BI \to \theta(x)$, torsion remains algebraic. After
elimination, the reduced action is~\cite{CalcagniMercuri2009,Mercuri2009}
\begin{equation}\label{eq:T1red}
  S_{\rm red} = S_{\rm EH} + S_{\rm Dirac}
  + \int\!\bigl[-\tfrac{1}{2}f(\theta)(\partial\theta)^2
    + \mathcal{L}_{\rm 4f}\bigr],
\end{equation}
where $\mathcal{L}_{\rm 4f} = -(3\pi G/2)
(\bar\psi\gamma^\mu\gamma^5\psi)^2$ is \textit{$\theta$-independent}.
Every term is reproducible by adding a standard ALP to GR. The Nieh-Yan
4-form is exact (total derivative) in 4D; its topological nature ensures
the gravitational origin washes out~\cite{CalcagniMercuri2009}. Cosmological
dynamics give $w = +1$ (stiff matter)~\cite{TaverasYunes2008}, the
opposite of dark energy. No mechanism generates $V(\theta) \sim \Lambda$.
Route~T1 closes.

\subsection{Route S1: parity CMB phenomenology}\label{sec:S1}

The minimal ECH action has no photon-polarization coupling. The
parity-odd operator $\varepsilon^{abcd}K_{ab}R_{cd}$ is purely
gravitational; the electromagnetic field strength does not appear in the
Holst term or the reduced action. Any assumed coupling produces
constant-$\beta$ birefringence with two free parameters for one
observable---identical to generic ALP phenomenology, with no distinctive
spin-torsion signature. Route~S1 closes.

\subsection{Summary}\label{sec:closure_summary}

\begin{table}[t]
\centering
\caption{Four tested routes and outcomes.}
\label{tab:closures}
\begin{ruledtabular}
\begin{tabular}{llll}
\textbf{Route} & \textbf{Target} & \textbf{Method} & \textbf{Failure} \\
\hline
Track B & Condensate & Fierz + NJL & $\Gsp < 0$ at $\BI < 1$ \\
G v1 & $\Gamma[e]$ & 1-loop & $\slashed{D}$ is $\BI$-free \\
T1 & $\theta(x)$ & Lit.\ review & $=$ standard ALP \\
S1 & Birefr. & Assessment & No photon coupling \\
\end{tabular}
\end{ruledtabular}
\end{table}


%% ============================================================
%%  6. STRUCTURAL LESSONS
%% ============================================================
\section{Structural Lessons}\label{sec:lessons}

The closures are not independent accidents. They reveal a hierarchy of
structural obstructions that progressively narrows the space of viable
geometric dark-energy models. Lessons~1--5 emerge from the minimal ECH
model and establish that algebraic torsion cannot carry information into
the IR. Lesson~6 emerges from the PGT extension and establishes that
even propagating torsion can fail---through a different mechanism---if
the mass and coupling share a common origin.

\subsection{Algebraic torsion washes out}

In the minimal ECH model, torsion is non-propagating. Its equation of
motion is algebraic: the contorsion $K$ is determined pointwise by the
fermion spin density. Torsion elimination is exact and complete---no
torsion-sector information survives into the reduced action except
through the four-fermion coupling constants. Any mechanism that requires
torsion to carry information into the infrared will fail in this model
class.

\subsection{Fixed $\BI$ leaves no IR fingerprint}

$\BI$ enters the reduced action only through $\Geff$ and $\Gsp$ in the
four-fermion interaction. These are UV-scale contact couplings.
At one loop, the fermion determinant is $\BI$-independent.
At the condensate level, the coupling is subcritical.
$\BI$ sets the UV-scale spin-torsion structure but does not generate
low-energy phenomena.

\subsection{Dynamical $\BI$ reduces to generic ALP}

The Nieh-Yan 4-form $N_4 = d(e^I \wedge T_I)$ is exact in 4D. Any
coupling of a pseudoscalar $\theta$ to $N_4$ reduces to boundary terms
plus standard kinetic/coupling structures after torsion elimination.
The gravitational origin is topologically protected from leaving a local
dynamical imprint. This is why Route~T1 collapses to a generic ALP.

\subsection{Parity alone is not a mechanism}

The ECH action's parity-odd sector is real, but without a derived
coupling to observable sectors (photons, matter), it produces no testable
prediction distinguishing the framework from generic alternatives. Parity
is a structural property of the theory, not by itself a mechanism or
signal source.

\subsection{No scale protection in the minimal model}

None of the tested routes addresses naturalness. The condensate route
produces exponentially suppressed scales but the wrong sign.
The one-loop route has standard renormalization.
The dynamical Immirzi route requires $V(\theta_0) = \Lobs$ by hand.
Any successor model must confront scale naturalness from the start.

\subsection{The mass-coupling lock}\label{sec:lesson6}

Even when torsion is promoted from algebraic to propagating (PGT), a new
obstruction emerges: the parameter that controls the torsion mass also
controls the torsion-matter coupling, and both go to zero together. In
the ghost-free $0^-$ axial torsion model, $m_B \propto 1/\sqrt{|t_3|}$
and $g_{\rm eff} \propto 1/\sqrt{|t_3|}$, so making the mode
cosmologically light automatically makes it cosmologically inert.

This is not a PGT-specific accident.
The canonical-normalization relation that produces this locking is
standard field theory; what is new is recognizing that it constitutes a
binding constraint on geometric dark-energy mechanisms---specifically,
that the ghost-free PGT torsion mode becomes cosmologically inert in
exactly the parameter regime required for dark energy.
The \textit{mass-coupling lock}
occurs generically when a single parameter controls both the kinetic
normalization and the mass of a geometric mode. Consider a field $\phi$
with Lagrangian $\mathcal{L} = -\tfrac{1}{2}Z(\partial\phi)^2
- \tfrac{1}{2}\mu^2\phi^2 + g\,\phi\,J$. Canonical normalization
$\phi_{\rm can} = \sqrt{Z}\,\phi$ gives $m = \mu/\sqrt{Z}$ and
$g_{\rm eff} = g/\sqrt{Z}$: mass and coupling are locked. In PGT,
$Z \propto |t_3|$ and $\mu^2 \propto \MPl^2$, producing exactly this
structure. The same pattern appears in Brans-Dicke theory, where the
scalar decouples as $\omega_{\rm BD} \to \infty$ (the GR limit
\textit{is} the decoupling limit), and in any geometric theory where the
propagating mode inherits its kinetic term from the gravitational action
without an independent mass-generating mechanism.

The lock is \textit{evaded} when mass and coupling are controlled by
independent parameters---for instance, through a Higgs-like mechanism
(separate VEV sets the mass), a gauge symmetry (coupling is fixed by a
charge), or an environmental mechanism (mass depends on background, not
on kinetic normalization). In the Standard Model, the $W$ boson mass
$m_W = gv/2$ is set by \textit{both} the gauge coupling $g$ and the
Higgs VEV $v$; taking $m_W \to 0$ by $v \to 0$ does not suppress $g$.
No analogous structure exists in standard PGT.

Any next-generation geometric dark-energy model must avoid this lock:
the mass and the coupling to observable sectors must be controlled by
independent parameters, or the light-mass limit must not simultaneously
decouple the field from matter.


%% ============================================================
%%  PART III: DIRECTIONS FORWARD
%% ============================================================

%% ============================================================
%%  7. REQUIREMENTS FOR NEXT-GENERATION FOUNDATIONS
%% ============================================================
\section{Requirements for Next-Generation Foundations}\label{sec:requirements}

The six structural lessons translate into necessary conditions that any
viable geometric dark-energy theory must satisfy simultaneously. These
are not aspirational guidelines; each corresponds to a demonstrated
failure mode. A candidate model that violates any single requirement is
excluded by the same mechanism that closed the corresponding route.

The first four requirements (DR1--DR4) were formulated from the
minimal-model closures; DR5 emerged from the Foundation~A (PGT)
assessment and addresses a structural obstruction not visible at the
minimal level.

\begin{description}[style=nextline,leftmargin=1em]
  \item[DR1: Survives reduction without washing out.]
    The model must contain at least one degree of freedom or structural
    feature that is present in both the full theory and the
    reduced/effective theory after all algebraic constraints are solved.
    \textit{Test:} write down the reduced action. Does the mechanism's key
    ingredient still appear?
    \textit{Closed by:} Lessons~1--3 (algebraic torsion, fixed $\BI$,
    Nieh-Yan topology).

  \item[DR2: Addresses tiny-scale naturalness explicitly.]
    The model must contain an explicit statement of why the relevant scale
    is small: shift symmetry, discrete symmetry, gauge symmetry, radiative
    stability, dynamical relaxation, or sequestering. ``The scale is a
    free parameter'' is not acceptable.
    \textit{Test:} if asked ``why $10^{-122}$?'', does the model have a
    structural answer?
    \textit{Closed by:} Lesson~5 (no scale protection in minimal model),
    Foundation~A Phase~2B (no symmetry protects PGT torsion mass).

  \item[DR3: Offers unique mechanism or distinctive observable.]
    The model must produce at least one prediction not reproducible by
    $\Lambda$CDM$+$generic ALP. If observationally identical to
    $\Lambda$CDM$+$ALP, the model has no scientific value regardless of
    theoretical motivation.
    \textit{Test:} name the distinctive prediction. Can a generic ALP
    reproduce it?
    \textit{Closed by:} Lesson~4 (parity alone is not a mechanism),
    Foundation~A Phase~2A (distinctive signatures suppressed by
    decoupling limit).

  \item[DR4: Fails cleanly if it doesn't work.]
    The model must be testable within a defined scope, with gates and kill
    criteria specified before computation. Negative results must be
    documentable.
    \textit{Test:} can you write the gate structure before calculating?

  \item[DR5: Avoids the mass-coupling lock.]
    The model's IR-relevant degree of freedom must not decouple from
    observable sectors in the light-mass limit. Specifically, the mass
    and the coupling to matter (or to photons, gravitational waves, or
    other observable channels) must be controlled by \textit{independent}
    parameters, or the light-mass limit must preserve the coupling at
    observable strength. A theory in which $m \to 0$ implies $g_{\rm eff}
    \to 0$ produces a cosmologically inert field regardless of its
    theoretical elegance.
    \textit{Test:} compute $g_{\rm eff}(m)$ in the canonically normalized
    theory. Does $g_{\rm eff}$ remain at gravitational strength or above
    as $m \to H_0$?
    \textit{Closed by:} Lesson~6 (mass-coupling lock in PGT).
\end{description}

\noindent
Foundation~A (PGT propagating torsion) passes DR1 and DR4 cleanly. It
fails DR2, DR3, and DR5. These failures are not independent: the
mass-coupling lock (DR5) is the structural mechanism by which DR3 fails
at cosmological parameter values. This pattern---a theory that is
healthy in the UV but observationally empty in the required IR
limit---is generic to theories where the mass and coupling share a
common kinetic-normalization origin (Sec.~\ref{sec:lesson6}). Evading
DR5 requires an independent mass-generating mechanism, not merely a
new choice of gauge group or field content.


%% ============================================================
%%  8. CANDIDATE FOUNDATION DIRECTIONS
%% ============================================================
\section{Candidate Foundation Directions}\label{sec:candidates}

Three directions were identified that each address at least one
structural lesson from Sec.~\ref{sec:lessons}. Foundation~A (propagating
torsion) has been tested through Phase~2 and found phenomenologically
empty (Sec.~\ref{sec:foundA}). Foundations~B and~C remain untested
starting points for future assessment using the methodology of
Sec.~\ref{sec:program}; both must now additionally satisfy DR5
(mass-coupling lock avoidance).

\subsection{Foundation A: Propagating torsion (tested)}\label{sec:foundA}

The most natural upgrade to the minimal ECH model is Poincar\'e gauge
theory (PGT), in which the spin connection is dynamical and torsion
propagates. This directly eliminates Lesson~1 (algebraic torsion wash-out)
by construction. We have conducted a systematic two-phase assessment of
this direction.

\textit{Phase~1 result.}
The quadratic PGT action (10 free parameters: 3 torsion-squared, 6
curvature-squared, plus $\Lambda_0$) admits three ghost-free single-mode
models, confirmed by independent analyses spanning four
decades~\cite{SezginNieuwenhuizen1980,Karananas2015}. The cleanest
candidate is Model~B: a $0^-$ pseudoscalar axial torsion mode, which
is ghost-free, parity-odd, and couples to the fermion axial current. Its
mass is $m_B = \MPl/(4\sqrt{\pi|t_3|})$, where $t_3$ is the
dimensionless torsion-squared coupling. A cosmologically light mass
($m_B \sim \text{meV}$) requires $|t_3|\sim 10^{58}$---the hierarchy is
transferred to a dimensionless coupling, not solved.

\textit{Phase~2 result: the mass-coupling lock.}
The critical finding concerns what happens in the cosmologically
interesting regime $|t_3| \gg 1$. The canonically normalized
matter coupling is
\begin{equation}\label{eq:geff}
  g_{\rm eff} \sim \frac{1}{\MPl\sqrt{|t_3|}}\,,
\end{equation}
so that $m_B$ and $g_{\rm eff}$ decrease \textit{together}: their ratio
$m_B/g_{\rm eff} \sim \MPl^2$ is fixed. At $|t_3|\sim 10^{58}$
(meV-scale mass), the matter coupling is $\sim\!10^{29}$ times weaker
than gravitational strength. The theory is perturbatively consistent (no
unitarity violation, no strong-coupling problem---the unitarity scale
\textit{rises} as $\sqrt{s_{\rm unit}}\sim \MPl\sqrt{|t_3|}$), but the
torsion mode cannot thermalize, cannot produce detectable fifth forces,
and cannot generate observable birefringence through its matter coupling.

Furthermore, no symmetry within PGT protects $m_B$: five candidate
symmetries (gauge, shift, chiral, topological, supersymmetric) were
investigated and excluded; the 't~Hooft naturalness criterion fails
($m_B = 0$ does not restore a symmetry); and graviton loops generate
$\delta m_B^2 \sim \MPl^2/(16\pi^2)$, requiring fine-tuning of
$1$~in~$10^{57}$ for $m_B \sim \text{meV}$.

\textit{Structural lesson.}
Large $|t_3|$ is a \textit{decoupling limit}, not a strong-coupling
limit. The same parameter that makes the mass cosmologically light also
kills the distinctive observational signatures. The reduced action is
GR~$+$~massive Proca pseudovector (structurally richer than GR~$+$~ALP,
with 3~DOF and axial-current coupling), but all distinctive features are
suppressed to irrelevance at cosmological parameter values. The mode
becomes indistinguishable from generic quintessence---only the Lagrangian
knows it is torsion.

\textit{DR assessment (post-test):}
DR1 (survives reduction): \textbf{pass}---the $0^-$ mode propagates and
is ghost-free.
DR2 (scale naturalness): \textbf{fail}---no symmetry protection, graviton
loop destabilizes mass.
DR3 (distinctive observable): \textbf{fail}---matter coupling too
suppressed at cosmological $|t_3|$.
DR4 (clean failure): \textbf{pass}---the decoupling limit is a clean,
physically interpretable negative result.

\textit{Verdict:} Foundation~A is structurally viable but
phenomenologically empty for dark energy. It resolves the algebraic
wash-out problem (Lesson~1) but encounters a new obstruction: the
\textit{mass-coupling lock} (Lesson~6, Sec.~\ref{sec:lesson6}).

\subsection{Foundation B: Symmetry-protected geometric pseudoscalar}
\label{sec:foundB}

The dynamical Immirzi field failed because the Nieh-Yan density is exact
in 4D. But in metric-affine gravity (where the connection is not
metric-compatible), the Nieh-Yan form is \textit{not} an exact
differential. A pseudoscalar coupled to this non-topological Nieh-Yan
density could retain geometric fingerprints after reduction, protected
by a shift symmetry that keeps its mass small.

\textit{Why it avoids old failures:}
ALP collapse (Lesson~3) is avoided if the non-topological coupling
retains geometric information. Scale protection (Lesson~5) is provided
by the shift symmetry.

\textit{Biggest risk:}
This may be Route~T1 in better clothing. The burden is entirely on
demonstrating that the metric-affine Nieh-Yan coupling is genuinely
non-topological and retains non-generic content after reduction.

\textit{First check:}
Is the Nieh-Yan density non-topological in metric-affine gravity? This
is a precise mathematical question with a checkable answer.

\textit{DR assessment:}
DR1: depends on the first-check answer. If the Nieh-Yan coupling reduces
to a total derivative after non-metricity is eliminated, this is
Route~T1 again.
DR2: shift symmetry provides natural scale protection.
DR3: if the non-topological coupling retains geometric content, the
pseudoscalar is non-generic (different from standard ALP).
DR4: yes---the topological/non-topological question is mathematically
definite.

\subsection{Foundation C: Vacuum sequestering}\label{sec:foundC}

Instead of generating a vacuum energy of the right magnitude, one can
attempt to cancel or sequester vacuum energy dynamically. The
Kaloper-Padilla mechanism enforces, via global constraints in the action,
that radiative corrections to the cosmological constant do not gravitate.

\textit{Why it avoids old failures:}
Scale naturalness (Lesson~5) is addressed directly. The small scale is
protected by a structural constraint, not generated by a local mechanism.

\textit{Biggest risk:}
Weinberg's no-go theorem~\cite{Weinberg1989} constrains local adjustment
mechanisms. Sequestering evades this by being global (non-local), but
this non-locality is controversial. Pure sequestering may set $\Lambda$
small but predict no distinctive observable (failing DR3).

\textit{First check:}
Write down the Kaloper-Padilla action~\cite{KaloPerPadilla2014}
in first-order EC gravity. Does the torsion sector modify the
sequestering constraint?

\textit{DR assessment:}
DR1: sequestering is a global mechanism, not reducible by algebraic
constraints---survives by construction.
DR2: directly addresses scale naturalness---this is its primary
motivation.
DR3: questionable---pure sequestering may set $\Lambda$ small but predict
no observable beyond $\Lambda$CDM. This is the critical weakness.
DR4: yes---whether torsion modifies the constraint is checkable.


%% ============================================================
%%  9. DISCUSSION
%% ============================================================
\section{Discussion}\label{sec:discussion}

\subsection{Current status of the framework}

The spin-torsion dark-energy framework occupies an honest intermediate
position. It is \textit{more} than generic $\Lambda$CDM: the scaling
ansatz~\eqref{eq:rhoL} provides a motivated parametric story for the
dark-energy scale, reducing fine-tuning by 115 orders of magnitude.
The parity-odd structure is real and could connect to observed
birefringence if a photon coupling is established.

It is \textit{less} than a first-principles derivation: $w = -1$ is
assumed, the scaling ansatz is not a mechanism, the observational
signatures are consistency checks rather than predictions, and
both the minimal model class and its most natural extension (PGT
propagating torsion) have been systematically tested and found
insufficient for distinct structural reasons.

\subsection{The closures as scientific contribution}

Negative results in theoretical physics are undervalued but scientifically
important. The five-closure program (four minimal routes plus the PGT
Foundation~A assessment):
\begin{itemize}[leftmargin=*]
  \item Establishes that both the minimal model and its most natural
    extension fail, for distinct structural reasons.
  \item Maps the failure mechanisms: algebraic torsion wash-out
    (minimal model) and the mass-coupling lock (PGT extension).
  \item Identifies five concrete requirements (DR1--DR5) for any
    successor, with DR5 (mass-coupling lock avoidance) being a
    previously unrecognized constraint.
  \item Demonstrates a reproducible methodology for disciplined model
    testing that can be applied to future candidates.
\end{itemize}

\subsection{Comparison with other programs}\label{sec:related}

The structural lessons identified here are not specific to our framework.
Any geometric dark-energy program based on minimal ECH gravity faces the
same algebraic-torsion wash-out. We position this work relative to five
classes of approaches.

\paragraph{Phenomenological dark energy.}
The $\Lambda$CDM$+\Delta\Neff$ extension analyzed here is statistically
equivalent to any additional relativistic species model (sterile
neutrinos, early dark energy). The framework adds theoretical
motivation for expecting $\Delta\Neff > 0$ (bounce particle production)
and a scaling argument for the $\Lambda$ scale, but the MCMC fits
themselves are standard. Quintessence and $w(z)$CDM models achieve
comparable or better tension reduction with different parameter
trade-offs; our advantage is the fine-tuning reduction from $10^{120}$
to $10^5$.

\paragraph{EC and Holst torsion cosmology.}
Pop\l{}awski~\cite{Poplawski2012} invokes torsion-induced four-fermion
condensates to source dark energy; our Track~B closure shows the S/P
channel is repulsive at $\BI < 1$, and the coupling is $175\times$
subcritical. Liu \etal~\cite{ECTorsionDESI2025} find EC torsion
preferred over $\Lambda$CDM by DESI~DR2 ($\Delta\text{AIC}\approx-6$),
but their torsion coupling is phenomenological---the same first-principles
gap applies. Legner, Handley \& Barker~\cite{Legner2025} proposed
Torsion Condensation (TorC) within Poincar\'e gauge theory, showing
that propagating torsion allows a larger inferred $H_0$---directly
supporting Foundation~A (Sec.~\ref{sec:foundA}). Fabbri~\cite{Fabbri2025}
showed that propagating torsion as an axial-vector field naturally
violates energy conditions, supporting the bounce mechanism.

\paragraph{Parity violation in gravity.}
Freidel \etal~\cite{Freidel2005} and Mercuri~\cite{Mercuri2006,Mercuri2009}
established the Holst term's parity-violating consequences and showed the
Barbero-Immirzi parameter is renormalized at one loop. Shapiro \&
Teixeira~\cite{ShapiroTeixeira2014} computed one-loop divergences with
external currents. These works establish the existence of a parity-odd
coefficient but do not extract a finite vacuum energy or trace it through
inflationary dilution. Our closure program (Branch~G) shows the strict
one-loop effective action is $\BI$-independent after torsion
elimination---a sharpening of the limitations implicit in these
earlier calculations.

\paragraph{Metric-affine and propagating torsion.}
In metric-affine gravity (non-metricity + torsion), the Nieh-Yan
density is \textit{not} an exact differential---potentially avoiding
the topological wash-out that closed Route~T1. In Poincar\'e gauge
theory (PGT), the spin connection is dynamical and torsion propagates
with massive spin-2 and spin-0 modes. The ghost-free parameter subspace
has been mapped by Sezgin \& van
Nieuwenhuizen~\cite{SezginNieuwenhuizen1980}, Yo-Nester, and
Karananas~\cite{Karananas2015}. Our Foundation~A assessment
(Sec.~\ref{sec:foundA}) confirms that ghost-free propagating modes
exist but reveals a mass-coupling lock: the same parameter that produces
a cosmologically light mass also suppresses the matter coupling to
irrelevance (Lesson~6). These approaches resolve Lesson~1 but encounter
a new structural obstruction.

\paragraph{Vacuum sequestering.}
The Kaloper-Padilla mechanism~\cite{KaloPerPadilla2014} enforces via
global constraints that radiative corrections to $\Lambda$ do not
gravitate. This evades Weinberg's no-go theorem~\cite{Weinberg1989}
through non-locality. Whether torsion modifies the sequestering
constraint in EC gravity is unexplored (Foundation~C first check).
Pure sequestering may set $\Lambda$ small but predicts no distinctive
observable, failing our DR3 requirement.

\subsection{What would change this assessment}

The framework's status would be upgraded if:
\begin{enumerate}
  \item A geometric theory is found in which the light-mass limit does
    \textit{not} decouple the field from matter---i.e., a theory that
    evades the mass-coupling lock identified in Foundation~A.
  \item A derived photon-torsion coupling is established, producing a
    non-generic birefringence prediction.
  \item Non-perturbative or higher-loop effects are shown to generate
    $\BI$-dependent vacuum structure beyond the tested orders.
  \item An external development (new symmetry, UV completion) identifies
    a protected light geometric degree of freedom with independent
    mass and coupling parameters.
\end{enumerate}
None of these outcomes is assumed or expected. The Foundation~A result
shows that simply upgrading from algebraic to propagating torsion is not
sufficient: the structural requirements are more demanding than
previously appreciated.


%% ============================================================
%%  10. CONCLUSION
%% ============================================================
\section{Conclusion}\label{sec:conclusions}

We have presented a phenomenological dark-energy framework rooted in
Einstein-Cartan-Holst gravity, reported the systematic closure of four
minimal first-principles routes to deriving its central assumption
($w = -1$ at late times), and assessed the most natural next-generation
extension---propagating torsion in Poincar\'e gauge theory---finding it
structurally viable but phenomenologically empty for dark energy.

The framework provides a motivated scaling ansatz reducing fine-tuning
from $10^{120}$ to $\sim\!10^5$, cosmological fits partially reducing
the $H_0$ and $S_8$ tensions, and a parity-odd structure qualitatively
consistent with observed birefringence (without a derived photon
coupling). These phenomenological results survive all closures intact.

The systematic model testing establishes:
\begin{itemize}[leftmargin=*]
  \item The minimal ECH model fails because algebraic torsion washes
    out, the Nieh-Yan density is topological, and no scale protection
    exists (Lessons~1--5).
  \item Propagating torsion in PGT resolves the algebraic wash-out but
    encounters a mass-coupling lock: the parameter that produces a
    cosmologically light mass simultaneously suppresses the matter
    coupling to irrelevance, and no symmetry protects the mass
    (Lesson~6).
  \item Six structural lessons and five decision rules (DR1--DR5)
    constrain the space of viable geometric dark-energy completions.
\end{itemize}

Any future geometric dark-energy model must not only survive reduction
and address scale naturalness, but must avoid the mass-coupling lock:
the field must remain cosmologically relevant in the light-mass limit.
This is a nontrivial requirement that eliminates the most natural
extension of the minimal framework.

The derivation of late-time $w = -1$ from geometry and the identification
of a distinctive observational signature remain open problems. The
contribution of this work is threefold: a well-tested phenomenological
model, a systematic closure of the five most natural derivation
strategies, and a concrete set of structural requirements---especially
the mass-coupling lock---that any future geometric dark-energy proposal
must address.


%% ============================================================
%%  ACKNOWLEDGMENTS
%% ============================================================
\begin{acknowledgments}
The author thanks the anonymous referees for constructive feedback.
Computational resources for MCMC analyses were provided by RunPod.
This work made use of Cobaya~\cite{Cobaya2021} and CAMB.
\end{acknowledgments}


%% ============================================================
%%  APPENDICES
%% ============================================================
\appendix

\section{Notation and Conventions}\label{app:notation}

Natural units $c=\hbar=1$ throughout, with $8\pi G = \MPl^{-2}$.
Metric signature $(-,+,+,+)$. Reduced Planck mass
$\MPl = (8\pi G)^{-1/2} = 2.435 \times 10^{18}$~GeV.
$\gamma^5 = i\gamma^0\gamma^1\gamma^2\gamma^3$ with
$(\gamma^5)^2 = 1$. The vacuum energy density is
$\rho_\Lambda = \Leff\,\MPl^2/(8\pi)$; when we write
$\rho_\Lambda \approx (2.3\;\text{meV})^4$ this refers to the energy
density. Greek indices $\mu,\nu,\ldots$ denote spacetime coordinates;
Latin indices $a,b,\ldots$ denote tangent-space (Lorentz) indices.
The Levi-Civita symbol $\varepsilon^{abcd}$ is totally antisymmetric
with $\varepsilon^{0123}=+1$.

Key symbols:
\begin{itemize}[leftmargin=*]
\item $\BI$---Barbero-Immirzi parameter ($0.274 \pm 0.020$)
\item $K_{ab}$---contorsion tensor
\item $T^a_{\;\;bc}$---torsion tensor
\item $\alpha/M$---parity-odd coefficient (mass dimension $-1$)
\item $[(\alpha/M)\,\MPl]\sim 10^{-2}$---dimensionless one-loop
  suppression factor
\item $\Dinf$---inflationary dilution factor ($\sim e^{-3N_{\rm tot}}$)
\item $\Xi\equiv[(\alpha/M)\,\MPl]\,\Dinf$---inflationary suppression
  factor; $\rho_\Lambda = \Xi\,\MPl^4$
\item $\Geff = (3\kappa^2/16)\,\BI^2/(\BI^2+1)$---effective four-fermion
  coupling
\item $\Gsp \propto (\BI^2-1)/(\BI^2+1)$---scalar/pseudoscalar channel
  coupling
\item $\Leff$---effective cosmological constant
\item $\rhocrit$---bounce critical density ($\simeq 0.27\,\rhoPl$)
\end{itemize}

Abbreviations: LQG (Loop Quantum Gravity), EC (Einstein-Cartan),
ECH (Einstein-Cartan-Holst), PGT (Poincar\'e Gauge Theory),
ALP (axion-like particle), NJL (Nambu-Jona-Lasinio),
MCMC (Markov Chain Monte Carlo), AIC/BIC (Akaike/Bayesian Information
Criterion).


\section{Parameter Summary}\label{app:params}

\begin{table*}[t]
\caption{\label{tab:params} Complete parameter classification. Status
categories: \textit{Fixed} (external input), \textit{Ansatz}
(phenomenological scaling), \textit{Assumed} (not derived),
\textit{Fit} (MCMC), \textit{Derived} (from other parameters).}
\begin{ruledtabular}
\begin{tabular}{llcll}
\textbf{Parameter} & \textbf{Description} & \textbf{Value}
  & \textbf{Status} & \textbf{Notes} \\
\hline
\multicolumn{5}{l}{\textit{Framework parameters}} \\
$\BI$ & Barbero-Immirzi & $0.274 \pm 0.020$ & Fixed & LQG area spectrum \\
$\alpha/M$ & Parity-odd coeff. & $\sim 10^{-33}\;\text{eV}^{-1}$
  & Fit & One-loop motivated \\
$N_{\rm tot}$ & Total $e$-folds & $\sim 92$ & Assumed
  & Fitted to $\rho_\Lambda$ \\
$\Dinf$ & Inflation dilution & $\sim e^{-3N_{\rm tot}}$
  & Derived & From $N_{\rm tot}$ \\
$\Xi$ & Infl.\ suppression & $\sim 10^{-123}$
  & Ansatz & $=[(\alpha/M)\MPl]\Dinf$ \\
$w$ & Equation of state & $-1$ & Assumed & Not derived \\
\hline
\multicolumn{5}{l}{\textit{Standard cosmological parameters (MCMC)}} \\
$H_0$ & Hubble constant & $69.2\pm 0.8$\,\textsuperscript{a}
  & Fit & km/s/Mpc \\
$\sigma_8$ & Clustering ampl. & $0.785\pm 0.016$\,\textsuperscript{a}
  & Fit & \\
$S_8$ & $\sigma_8(\Omega_m/0.3)^{0.5}$ & $0.80\pm 0.02$
  & Derived & \\
$\Omega_m$ & Matter density & $0.310 \pm 0.008$ & Fit & \\
$\Omega_b h^2$ & Baryon density & $0.02237 \pm 0.00015$ & Fit & \\
$n_s$ & Spectral index & $0.9649 \pm 0.0042$ & Fit & \\
$\tau$ & Optical depth & $0.054 \pm 0.007$ & Fit & \\
\hline
\multicolumn{5}{l}{\textit{Extended parameters}} \\
$\Delta\Neff$ & Extra species & $\sim 0$\,\textsuperscript{a}
  & Fit & Phenomenological \\
$\Omega_k$ & Curvature & 0 & Fixed & Mandated by $N_{\rm tot}$ \\
$(\omega/H)_0$ & Vorticity & $< 5\times 10^{-11}$
  & Bounded & CMB isotropy \\
\end{tabular}
\end{ruledtabular}
\end{table*}

\textsuperscript{a}Original values from the tension dataset (includes
SH0ES $H_0$ prior). Independent verification
(Table~\ref{tab:verification}) finds $H_0 = 67.68 \pm 1.06$,
$\sigma_8 = 0.803 \pm 0.008$, $\Delta\Neff = -0.020 \pm 0.169$
(full-tension) and $H_0 = 67.79 \pm 1.09$,
$\sigma_8 = 0.812 \pm 0.009$, $\Delta\Neff = +0.065 \pm 0.17$
(Planck+BAO+SN).

Key parameter correlations: $r(H_0,\sigma_8)=-0.45$;
$r(\alpha/M,\Dinf)=-0.89$ (only their product $\Xi$ is
well-constrained). The strong anticorrelation means the framework
predicts $\rho_\Lambda = \Xi\,\MPl^4$ as a single effective
parameter, not two independent ones.


\section{Claims Classification}\label{app:claims}

\begin{table*}[t]
\caption{Claims classification: what is derived, what is assumed, what is
fit, and what has been retired by the closure program. This table
supersedes the classification in~\cite{Golden2026v1}.}
\label{tab:claims}
\begin{ruledtabular}
\begin{tabular}{llll}
\textbf{Claim} & \textbf{Status} & \textbf{Basis} & \textbf{Notes} \\
\hline
\multicolumn{4}{l}{\textit{Derived (structural results)}} \\
ECH action structure & Derived & Standard EC + Holst & Background \\
Torsion is algebraic (minimal ECH) & Derived & EC field equations & Structural \\
Perfect-square four-fermion & Derived & Torsion elimination & Exact result \\
$\Gsp < 0$ at $\BI < 1$ & Derived & Fierz rearrangement & Track B closure \\
One-loop $\slashed{D}$ is $\BI$-free & Derived & Torsion elimination & G v1 closure \\
Dyn.\ Immirzi $=$ ALP & Derived (lit.) & Calcagni-Mercuri 2009 & T1 closure \\
No photon coupling in min.\ ECH & Derived & Action structure & S1 closure \\
Nieh-Yan is topological in 4D & Derived (lit.) & Differential geometry & T1 closure \\
\hline
\multicolumn{4}{l}{\textit{Phenomenological (ansatz, fit, or assumption)}} \\
$\rho_\Lambda = \Xi\,\MPl^4$ & Ansatz & Scaling argument & Not a derivation \\
$w = -1$ at late times & Assumed & Observational input & Not derived \\
$\Delta\Neff$ from bounce & Phenomenological & Qualitative motivation & Consistent with 0 \\
$H_0$, $\sigma_8$, $S_8$ values & Fit & MCMC & Standard analysis \\
Fine-tuning $10^{120} \to 10^5$ & Parametric obs. & Sensitivity scan & Valid \\
Birefringence consistency & Consistency check & Assumed coupling & Not a prediction \\
Galaxy spin $A_0 \sim 0.003$ & Empirical & Hierarchical Bayes fit & Not derived \\
\hline
\multicolumn{4}{l}{\textit{Retired (closed by derivation/assessment program)}} \\
Birefringence as prediction & \textbf{Retired} & S1 closure & No photon coupling \\
$\BI$ drives late-time physics & \textbf{Retired} & B, G, T1 closures & UV-scale only \\
Condensate mechanism & \textbf{Retired} & Track B closure & $\Gsp < 0$ \\
$\Delta\Neff$ as distinctive signal & \textbf{Retired} & Verification & Generic to any $\Neff$ ext. \\
Galaxy spin as prediction & \textbf{Retired} & Gap analysis & $A_0$ not derivable \\
PGT torsion as DE mechanism & \textbf{Retired} & Foundation A & Mass-coupling lock \\
\end{tabular}
\end{ruledtabular}
\end{table*}


%% ============================================================
%%  REFERENCES
%% ============================================================

\begin{thebibliography}{99}

\bibitem{Weinberg1989}
S.~Weinberg,
Rev.\ Mod.\ Phys.\ \textbf{61}, 1 (1989).

\bibitem{Planck2018params}
Planck Collaboration,
Astron.\ Astrophys.\ \textbf{641}, A6 (2020).

\bibitem{Riess2022}
A.~G.~Riess \etal,
Astrophys.\ J.\ Lett.\ \textbf{934}, L7 (2022).

\bibitem{KiDS2021}
KiDS Collaboration,
Astron.\ Astrophys.\ \textbf{645}, A104 (2021).

\bibitem{DES2024}
DES Collaboration,
Phys.\ Rev.\ D~\textbf{105}, 023520 (2022).

\bibitem{DESI2024}
DESI Collaboration,
arXiv:2404.03002 (2024).

\bibitem{DESI2025DR2}
DESI Collaboration,
arXiv:2503.14738 (2025).

\bibitem{Ashtekar2011}
A.~Ashtekar and P.~Singh,
Class.\ Quant.\ Grav.\ \textbf{28}, 213001 (2011).

\bibitem{Freidel2005}
L.~Freidel, D.~Minic, and T.~Takeuchi,
Phys.\ Rev.\ D~\textbf{72}, 104002 (2005).

\bibitem{Mercuri2006}
R.~Mercuri,
Phys.\ Rev.\ D~\textbf{73}, 084016 (2006).

\bibitem{Mercuri2009}
R.~Mercuri,
Phys.\ Rev.\ Lett.\ \textbf{103}, 081302 (2009).

\bibitem{Poplawski2012}
N.~J.~Pop\l{}awski,
Phys.\ Rev.\ D~\textbf{85}, 107502 (2012).

\bibitem{ShapiroTeixeira2014}
I.~L.~Shapiro and P.~M.~Teixeira,
Class.\ Quant.\ Grav.\ \textbf{31}, 185002 (2014).

\bibitem{Shapiro2002}
I.~L.~Shapiro,
Phys.\ Rept.\ \textbf{357}, 113 (2002).

\bibitem{Vassilevich2003}
D.~V.~Vassilevich,
Phys.\ Rept.\ \textbf{388}, 279 (2003).

\bibitem{CalcagniMercuri2009}
G.~Calcagni and R.~Mercuri,
Phys.\ Rev.\ D~\textbf{79}, 104003 (2009).

\bibitem{TaverasYunes2008}
V.~Taveras and N.~Yunes,
Phys.\ Rev.\ D~\textbf{78}, 064070 (2008).

\bibitem{Eskilt2022}
J.~R.~Eskilt and E.~Komatsu,
Phys.\ Rev.\ D~\textbf{106}, 063503 (2022).

\bibitem{DiegoPalazuelos2025}
P.~Diego-Palazuelos and E.~Komatsu,
(2025).

\bibitem{Golden2026v1}
H.~Golden,
``Geometric Dark Energy from Spin-Torsion Cosmology:
Phenomenological Constraints and Correlated Signatures,''
Preprint HUBIFY-2026-001, v1.6 (2026).

\bibitem{Golden2026supplement}
H.~Golden,
``Systematic Closure of Minimal First-Principles Routes to Dark Energy
in Einstein-Cartan-Holst Gravity,''
Companion technical note (2026).

\bibitem{Cobaya2021}
J.~Torrado and A.~Lewis,
JCAP \textbf{05}, 057 (2021).

\bibitem{ECTorsionDESI2025}
J.~Liu, R.~Niu, and R.-G.~Cai,
arXiv:2503.02599 (2025).

\bibitem{Saadeh2016}
D.~Saadeh \etal,
Phys.\ Rev.\ Lett.\ \textbf{117}, 131302 (2016).

\bibitem{Legner2025}
C.~Legner, W.~J.~Handley, W.~E.~V.~Barker, and S.~Sherrill Ormondroyd,
arXiv:2507.09228 (2025).

\bibitem{Fabbri2025}
L.~Fabbri,
arXiv:2502.17979 (2025).

\bibitem{KaloPerPadilla2014}
N.~Kaloper and A.~Padilla,
Phys.\ Rev.\ Lett.\ \textbf{112}, 091304 (2014).

\bibitem{SezginNieuwenhuizen1980}
E.~Sezgin and P.~van Nieuwenhuizen,
Phys.\ Rev.\ D~\textbf{21}, 3269 (1980).

\bibitem{Karananas2015}
G.~K.~Karananas,
Class.\ Quant.\ Grav.\ \textbf{32}, 055012 (2015).

\end{thebibliography}

\end{document}
